\title{Large Language Models Can Automatically Engineer Features for Few-Shot Tabular Learning}

\begin{abstract}
Large Language Models (LLMs) have demonstrated notable capacity for addressing intricate and previously unseen reasoning challenges, making them highly promising for tabular learning—a core requirement across numerous practical fields. This work introduces \textsf{FeatLLM}, an innovative in-context learning methodology that leverages LLMs as automated feature creators, producing an enhanced input dataset tailored for tabular prediction tasks. The new features constructed by the LLM enable reliable estimation of class probabilities via a straightforward downstream machine learning algorithm such as linear regression, facilitating highly effective few-shot learning. Uniquely, \textsf{FeatLLM} limits inference to the use of the generated features within this basic predictive model, without engaging the LLM for every test instance. In contrast to existing LLM-based models, \textsf{FeatLLM} eradicates the requirement of querying the LLM for each data point at inference, functions with only API-based LLM access, and circumvents prompt length constraints. Comprehensive evaluation on diverse tabular benchmarks demonstrates that \textsf{FeatLLM} produces superior rule sets, achieving considerable improvements (an average gain of 10\%) over competitors such as TabLLM and STUNT.
\end{abstract}

\section{Introduction}
Large Language Models (LLMs) have recently established themselves as capable systems for generating contextually appropriate responses to tasks they were not explicitly trained for, with no need for fine-tuning on task-specific data~\cite{creswell2022selection,imani2023mathprompter,taylor2022galactica}. These models are pretrained on large-scale textual datasets, imbuing them with a vast repository of implicit world knowledge, which can obviate the necessity of hand-crafted domain expertise for a variety of tasks~\cite{Genomes2021,singhal2023large,wilcox2003role}. This has been pivotal in advancing the automation of numerous applications, such as dialogue processing~\cite{finch2023leveraging} and program debugging~\cite{fan2023automated}. By incorporating concise task specifications and illustrative examples into their prompts, LLMs demonstrate the ability to interpret contextual information, focus on salient attributes, and relevant training samples~\cite{brown2020language,zhao2023survey}. Such behavior suggests that LLMs can function effectively as feature engineers, systematically extracting and highlighting relevant information from data. 

Recently, researchers have harnessed the reasoning skills and background knowledge of LLMs for the purpose of tabular learning. For example, LLMs can utilize concepts like the association between advanced age and higher prevalence of certain diseases to improve disease risk estimation tasks. In these efforts, tabular information and relevant feature descriptions are first expressed in natural language and serialized before being supplied to the LLM~\cite{dinh2022lift,wang2023anypredict}. These approaches typically employ either in-context learning~\cite{nam2023semi} or efficient adaptation methods that adjust only a subset of parameters~\cite{dinh2022lift,hegselmann2023tabllm}. Leveraging the prior knowledge encoded in LLMs has been shown to consistently surpass traditional tabular modeling techniques, especially in settings characterized by limited training data~\cite{hegselmann2023tabllm}.

Existing tabular learning strategies leveraging LLMs exhibit several shortcomings. For end-to-end predictive tasks, these methods require conducting an LLM inference for each data instance, leading to substantial computational demands. Additionally, achieving strong predictive performance generally necessitates fine-tuning the LLM, which presents practical obstacles when utilizing models that lack openly available training infrastructure or fine-tuning APIs. This issue has grown in significance as high-performing, recently introduced LLMs typically provide only constrained inference access through APIs~\cite{anil2023palm,openai2023gpt4}. Another notable limitation is that most approaches do not handle long prompts well~\cite{liu2023lost}: as the dimensionality of tabular datasets increases, resulting in lengthier serialized input texts, these models either fail to operate or suffer noticeable performance deterioration. Collectively, such barriers hinder the adoption of LLM-based tabular models for practical, real-world problems.

Our methodology reorients the application of LLMs in this domain: rather than using them exclusively for generating end-to-end predictions, we harness their capabilities for the purpose of feature construction. Specifically, we introduce \textsf{FeatLLM}, a new framework centered around in-context learning that seeks to elicit and interpret the implicit “criteria” that drive LLM decision-making. As an illustrative example, in disease classification, the LLM can analyze input features and output explicit diagnostic rules delineating the feature patterns corresponding to the disease state. These induced rules, or criteria, serve as blueprints to derive transformed features that substitute or augment the original inputs for downstream modeling tasks. This reallocation of the LLM’s utility—transitioning from a direct predictor to a feature generator—facilitates the deployment of simpler, more efficient predictive models, thus minimizing inference delays. Taking advantage of in-context learning, features can be obtained by merely providing demonstration examples in the prompts, obviating the need for conventional model training. To further address scalability issues related to prompt length, ensemble techniques such as feature bagging~\cite{breiman2001random} can be integrated to curtail prompt size and enhance overall robustness.

\textsf{FeatLLM} divides the feature engineering process into two principal stages: (1) comprehending the given problem and deducing how features are related to the predictive targets; and (2) leveraging insights gained from the first stage, together with few-shot demonstration examples, to construct discriminative rules for class separation. This organized reasoning sequence encourages LLMs to disregard irrelevant features while deriving candidate rules. The resulting set of rules is implemented on the dataset—by parsing them as programmatic instructions—yielding transformed binary features that indicate whether individual samples satisfy each rule. Subsequently, a simple model is trained using these features to produce class probability estimates. To enhance robustness, an ensemble method is employed, where the generation of features is repeated under a feature bagging strategy. Tuning the numbers of sampled features and examples in this procedure alleviates the issue of large prompt size. Free from any training phase and incurring minimal inference overhead, \textsf{FeatLLM} is thus positioned as a practical tool for a broad range of applications.

Our assessment of \textsf{FeatLLM} covers 13 distinct tabular datasets in low-data contexts, demonstrating consistently strong and stable results. The method capitalizes on both pre-existing LLM knowledge and knowledge extracted from data to produce effective features. Our approach surpasses leading few-shot learning benchmarks under various experimental conditions. The associated implementation is openly provided through an anonymized GitHub repository at~\texttt{https://github.com/Sungwon-Han/FeatLLM}.

\section{Related Work}

\subsection{Few-Shot Learning with Tabular Data} 

Progress in few-shot learning has made it possible to derive broadly useful data representations even when annotation is scarce~\cite{chen2018closer}. Although initial research concentrated on image data, the adaptability of few-shot learning now extends to several data types, such as text and audio, with growing application to tabular datasets~\cite{majumder2022few,nam2023stunt,schick2021exploiting}. The challenge in such scenarios lies in the propensity of models to pick up on irrelevant or spurious correlations due to limited supervision~\cite{bartlett2021deep}. As a result, contemporary research has focused on integrating external information sources or encoding domain knowledge to introduce effective inductive biases in the learning process. For example, semi-supervised learning makes use of vast amounts of unlabeled data, augmented appropriately, to bolster labeled samples~\cite{ucar2021subtab,yoon2020vime}, while unsupervised pre-training methods aim to provide a solid initialization irrespective of the task~\cite{somepalli2021saint}. Approaches within this paradigm typically involve reconstructive objectives—such as reconstructing masked or corrupted inputs~\cite{arik2021tabnet,majmundar2022met}—or utilize augmented data in contrastive frameworks~\cite{bahri2022scarf,chen2020simple}. Nevertheless, designing effective augmentations for tabular data is complicated by its inherent heterogeneity, often resulting in suboptimal performance for few-shot learning applications~\cite{nam2023stunt}.

A more recent trend involves training models over a large and diverse array of tabular datasets, regardless of their relation to the downstream target task, and then adapting these models via transfer learning~\cite{zhang2023generative,levin2022transfer,zhu2023xtab}. TransTab~\cite{wang2022transtab}, for example, employs transformer architectures that not only use table values but also leverage semantic information from column names to discern inter-column relationships. This allows for the transfer of accumulated tabular knowledge to new tasks, thereby enabling zero-shot and few-shot inference. In another instance, TabPFN~\cite{hollmann2022tabpfn} generates synthetic data by combining various distributions to pre-train a Bayesian neural network. Post pre-training, the model requires no further updating, handling both training and testing samples of the target task jointly during inference. Knowledge obtained from diverse datasets proved to outperform classical tree-based algorithms and demonstrated clear advantages in low-shot learning scenarios.

\subsection{Language-Interfaced Tabular Learning}

An alternative and rapidly developing direction to infuse prior knowledge into tabular models is through the use of large language models (LLMs)~\cite{anil2023palm,openai2023gpt4}. Owing to pre-training on massive, varied textual resources, LLMs excel at generalization, even for novel tasks, showcasing effectiveness across multiple application areas~\cite{brown2020language}. Recent investigations have sought to apply this expansive prior knowledge within tabular learning, generally by transforming tabular inputs into textual formats, which are then processed as input prompts by LLMs~\cite{dinh2022lift,hegselmann2023tabllm,wang2023anypredict}. Prompt construction typically includes not only the raw data, but also feature and task descriptors, enabling LLMs to reason about relationships among variables to produce meaningful predictions. Advances in this vein have either employed lightweight, parameter-efficient fine-tuning methods such as LoRA~\cite{hu2021lora} and IA3~\cite{liu2022few}, or opted for in-context learning whereby a handful of task examples are introduced in the prompt without further training~\cite{dinh2022lift,hegselmann2023tabllm,nam2023semi,slack2023tablet}.

Although the previously described approaches demonstrate substantial success in enhancing few-shot tabular learning, their reliance on routing every inference through the LLM introduces notable obstacles, particularly for practical deployment in industry settings. This research seeks to move away from the typical black-box methodology where every prediction is processed end-to-end by an LLM. Rather, it proposes to utilize the LLM specifically for the direct extraction of the governing conditions or decision rules for each class. \textsf{FeatLLM} interprets the rules identified by the LLM as executable code, which is then employed to construct novel features. This enables the system to produce predictions independent of the LLM at inference time, and takes advantage of in-context learning to derive rules by leveraging prior domain knowledge as well as the small number of labeled examples.

\section{Methods}
The goal of \textsf{FeatLLM} is to extract "rules"—namely, the set of conditions defining each class label—from a limited collection of samples, as illustrated in Figure~\ref{fig:main_model}. Once these rules are output by the LLM, they are programmatically parsed and transformed into new binary features. For each given data point, these features indicate which of the class-specific rules are satisfied. Subsequently, these binary features are aggregated through a dot product with non-negative trainable weights, providing a class likelihood estimate. By fitting this model, we can assess the relative significance of each feature and its contribution to class membership (see Section~\ref{Sec:training}). This learning procedure is iteratively executed with bagging, and results from these repetitions are merged using ensemble techniques. Further elaboration of the problem formulation and methodology for each step are provided in the subsections below.

\vspace*{-0.12in}
\paragraph{Problem formulation.} Consider a tabular dataset containing $N$ labeled examples, denoted as $\mathcal{D} = \{(\mathbf{x}^i, \mathbf{y}^i)\}_{i=1}^{N}$. In this collection, each sample consists of a $d$-dimensional input vector $\mathbf{x}^i$, where each dimension corresponds to one of the $d$ input features, and feature names are captured by $F = \{f_j\}_{j=1}^{d}$. In addition to names, separate feature definitions are supplied in natural language. In the context of classification, the labels $\mathbf{y}^i$ take values from a fixed set of class categories. For $k$-shot learning tasks, the model is trained using a randomly selected subset comprising $k$ examples, where $k < N$.

\subsection{Prompt Design for Extracting Rules}
\label{Sec:prompt_design}
To foster more precise rule extraction from the LLM, we structure the problem-solving prompt to align with the step-by-step reasoning often used by domain experts when tackling tabular data problems. The designed prompt includes three core segments (see Fig.~\ref{fig:prompt_for_rule} and further details in Appendix~\ref{sec:example_rule}):

\vspace*{-0.12in}
\paragraph{Basic information description.} This section supplies the fundamental context needed for task resolution (highlighted in orange in Fig.~\ref{fig:prompt_for_rule}), covering the following: the formulation of the problem with an explicit question and label definition, feature descriptions, and a handful of example cases. The task overview is framed as an explicit question (for instance, ``Does this patient's myocardial infarction complications data indicate chronic heart failure? Yes or no?"). Further illustration of task queries can be found in Appendix~\ref{sec:task_desc}. Feature information details its datatype and, where suitable, its definition; for categorical inputs, potential categories may be specified. Example cases are presented by serializing a small selection of training instances into text format, including the correct label for each. This serialization applies a format in line with prior works~\cite{dinh2022lift,hegselmann2023tabllm}:

\begin{align}
&\text{Serialize}(\mathbf{x}^i,\mathbf{y}^i, F) = \nonumber \\ 
&``\ f_1\ \text{is}\ \mathbf{x}^i_1.\ ... \ f_d\  \text{is}\  \mathbf{x}^i_d.\ \ \text{Answer:}\  \mathbf{y}^i\ ”
\end{align}

\vspace*{-0.12in}
\paragraph{Reasoning instruction.} Our goal is to improve the LLM's reasoning capabilities by offering targeted reasoning instructions (highlighted in blue in Fig.~\ref{fig:prompt_for_rule}). The reasoning instruction commences with an introductory statement, drawing inspiration from the chain-of-thought methodology~\cite{wei2022chain}. Subsequently, we structure the LLM’s reasoning workflow into two discrete phases. In the initial phase, the LLM is prompted to deduce potential causal links or associations between features and the specific task, intentionally without relying on example demonstrations and instead leveraging general domain knowledge or common sense. This design encourages the LLM to utilize its prior conceptual knowledge about the context. The subsequent phase directs the LLM to synthesize information from both the earlier inference and provided example demonstrations to establish class-specific rules. Segmenting the reasoning process in this manner helps the LLM avoid the pitfall of attending to insignificant or misleading correlations among unrelated features, thereby promoting a focus on more pertinent attributes. To balance between insufficient and excessive rule extraction—which can result in either underfitting or an overly cumbersome prompt—we constrain the process by fixing the number of rules to be extracted per class to 10\footnote{Section~\ref{sec:analysis} provides an analysis regarding rule quantity selection.}. Furthermore, when articulating these rules, we instruct the LLM to consult the predefined types of each feature from the basic information segment to circumvent the risk of mismatched rule-feature types.

\vspace*{-0.12in}
\paragraph{Response instruction.} To streamline downstream parsing and utilization of the LLM-generated rules, we provide explicit instructions guiding the LLM’s formatting of its outputs (as indicated by the yellow text in Fig.~\ref{fig:prompt_for_rule}). The prompt delineates the expected response layout for each class (i.e., the response format) as well as the structural template that each individual rule should adhere to (i.e., the rule format). These directives equip the LLM to choose formatting strategies that correspond to the feature types in question. Should further specification be omitted, the LLM retains the flexibility to employ logical connectors such as AND or OR to combine various rules. An illustration of such output is available in Appendix~\ref{sec:outcome-rule}.

\subsection{Inferring Class Likelihood via Rules}
\label{Sec:training}

\paragraph{Parsing rules for feature generation.} The rules derived from the prior step are employed to engineer new binary features for each class label. Each feature denotes whether a data sample adheres to the respective class rule (coded as '0' or '1'), which serves as an indicator of the sample's likelihood of class membership. Since these rules originate from the LLM in natural language form, an effective parsing mechanism is crucial to facilitate the automatic creation of these features. Despite our efforts to standardize rule format and response structure for smoother parsing, LLM-generated text may still exhibit inconsistencies or unexpected syntactic variations~\cite{kovavc2023large}. For example, the LLM might rephrase symbols such as ``$>$” or ``$<$” with textual equivalents like ``greater than" or ``less than," introducing additional complexity into the parsing process.

In order to mitigate the difficulties posed by such noisy outputs, rather than implementing intricate parsing algorithms, we capitalize on the LLM’s own semantic processing strengths. The LLM’s proficiency at discerning intended meaning, even with syntactic deviations, together with its capacity for generating functional code (due to training on source code data), can be harnessed for this task. We construct prompts that include the function signature, explanations of inputs and outputs, and the inferred rules, then submit these prompts to the LLM. Illustrative prompts and outputs are provided in Figures~\ref{fig:prompt_for_parsing} and ~\ref{sec:example_parsing}-\ref{sec:example_parse_outcome} in the Appendix. The code produced by the LLM is subsequently run using Python’s \texttt{exec()} to translate these parsed rules into features within the dataset.

\vspace*{-0.12in}
\paragraph{Inferring class likelihood.} After extracting the binary features from rules corresponding to each class, an initial approach to evaluate the likelihood that a sample belongs to a particular class is to tally the number of satisfied rules associated with that class—equivalently, this is the sum of the binary features for the class. Nonetheless, rules and their corresponding features may hold varying degrees of importance, making it necessary to quantify this differential relevance through learning on training data. To accomplish this, we propose to fit a standard machine learning (ML) model, applying it separately to each set of class-specific binary features. As an instance, a bias-free linear model can be used. Let us define the binary feature vector for sample $\mathbf{x}^i$ and class $k$ as $\mathbf{z}_k^i \in \{0, 1\}^R$, with $R$ representing the number of rules per class, and $c$ indicating the total number of classes. The class probabilities $\mathbf{p}^i$ are estimated as:

\begin{align}
\text{logit}_k^i &= \max(\mathbf{w}_k, 0) \cdot \mathbf{z}_k^i, \nonumber \\
\mathbf{p}^i &= \text{Softmax}({[\text{logit}_{1}^i, ..., \text{logit}_{c}^i]}). \label{eq:infer_prob}
\end{align}

In this formulation, $\mathbf{w}_k$ stands for the trainable parameters in the linear model and encapsulates the relative importance of each feature. By constraining these weights to be non-negative, we guarantee that none of the features from the LLM acts to decrease the predicted likelihood of any class during model training. After computing the logit vector, the softmax operation transforms it into normalized class probabilities. Training proceeds by minimizing the cross-entropy loss with respect to the weights $\mathbf{w}_k$, utilizing a limited set of labeled examples.

\vspace*{-0.12in}
\paragraph{Ensembling with bagging.} To further enhance robustness, the above procedure is iterated multiple times to generate a set of inference models whose outputs are ultimately aggregated through ensembling. Denote the ensemble of learned weights from $T$ independent runs as $\{\mathbf{w}[t]\}_{t=1}^T$; the predictive output is obtained by averaging the per-class probabilities $\mathbf{p}^i$—calculated for each instance of $\mathbf{w}[t]$ in accordance with Eq.~\ref{eq:infer_prob}—across all $T$ models.

To incorporate stochasticity into the rule generation process for ensembles, we set a relatively high temperature during LLM inference. Additionally, we shuffle the sequence of few-shot examples, motivated by evidence that the ordering of demonstrations influences the response generated by the LLM~\cite{lu2022fantastically}\footnote{Further examination of rule diversity is available in Appendix~\ref{sec:diversity}}. In order to leverage this diversity and enhance prediction robustness, we employ a bagging strategy whereby, for each ensemble iteration, a random subset of features or data instances is selected. This method becomes particularly beneficial as the dataset or feature space grows in size. The ensemble framework thus brings two significant benefits. Firstly, inaccuracies from faulty rules in individual trials may be offset by correct rules identified in other ensemble members, thereby strengthening overall model resilience. This behavior reflects the LLM’s intrinsic self-consistency~\cite{wang2022self}, as rules reproduced across multiple samples tend to have greater validity. Secondly, this ensemble mechanism mitigates the constraint imposed by the LLM’s limited prompt capacity. When tabular datasets exceed this size, bagging decreases the input demand, preserving robust model performance. While an individual rule set may not encompass the dependencies of every feature, aggregating several weak learners generated in separate ensemble rounds compensates for possible feature or sample omissions in any single trial.

\section{Experiments}
We test the capabilities of \textsf{FeatLLM} across a range of tabular datasets and perform comprehensive analyses to elucidate various facets of our approach. Owing to space limits, supplementary experiments—including experiments with alternative LLM architectures, in-depth qualitative studies, and more—are detailed in the Appendix.

\subsection{Performance Evaluation}
\paragraph{Datasets.}
Our study involves 13 datasets addressing both binary and multi-class classification challenges: (1) Adult~\cite{asuncion2007uci}, where the goal is to predict if a person's annual income exceeds \$50,000; (2) Bank~\cite{moro2014data}, which predicts customer subscription to a term deposit; (3) Blood~\cite{yeh2009knowledge}, focused on forecasting repeat blood donation behavior; (4) Car~\cite{kadra2021well}, tasked with determining the assessed quality of an automobile; (5) Communities~\cite{misc_communities_and_crime_183}, aiming to estimate crime rates within specific regions; (6) Credit-g~\cite{kadra2021well}, which forecasts an individual’s likelihood of being a favorable or unfavorable credit risk; (7) Diabetes\footnote{\texttt{kaggle.com/datasets/uciml/pima-indians-diabetes-database}}, used for identifying diabetes diagnoses; (8) Heart\footnote{\texttt{kaggle.com/datasets/fedesoriano/heart-failure-prediction}}, for predicting incidents of coronary artery disease; and (9) Myocardial~\cite{misc_myocardial_infarction_complications_579}, which ascertains whether a person has chronic heart failure.

Additionally, we incorporate several newer and synthetic datasets that have received limited attention in prior work and were not included during the LLM pre-training phase: (10) Cultivars, used to predict the projected grain yield of specific soybean varieties\footnote{\texttt{archive.ics.uci.edu/dataset/913}}; (11) NHANES, for determining an individual's age group based on their record\footnote{\texttt{archive.ics.uci.edu/dataset/887}}; (12) Sequence-type, which involves classifying a provided sequence as arithmetic, geometric, Fibonacci, or Collatz; and (13) Solution-mix, in which the task is to predict whether the concentration percentage of a four-solution mixture is greater than 0.5. The first two of these datasets were only recently submitted to the UCI Machine Learning Repository after September 2023, thereby excluding them from the pre-training regime of the LLM API (gpt-3.5-turbo-0613) employed by our system. The other two are synthetic datasets, ensuring that the LLM API could not have memorized their contents, which makes them suitable for evaluation purposes.

The chosen datasets present diversity in both their scale and level of complexity, as outlined in Table~\ref{Tab:basic-desc}. Each one is accompanied by explicit attribute names and explanations. Datasets such as `Communities' and `Myocardial' contain more than 100 features, leading to difficulties due to the LLM’s restricted context window.

\vspace*{-0.12in}
\paragraph{Baselines.}
We evaluate our approach against nine baseline methods. The initial three represent classic tabular data algorithms: (1) Logistic regression (LogReg), (2) XGBoost~\cite{chen2016xgboost}, and (3) RandomForest~\cite{ho1995random}. The subsequent two baselines, (4) SCARF~\cite{bahri2022scarf} and (5) STUNT~\cite{nam2023stunt}, make use of unlabeled data; however, in many real-world contexts, acquiring vast amounts of unlabeled data may not be practical. An additional recently proposed baseline is (6) TabPFN~\cite{hollmann2022tabpfn}, which leverages a broad array of synthetically generated data distributions to pre-train its model. The remaining three baselines involve LLM-based techniques: (7) In-context learning (In-context)~\cite{wei2022emergent}, (8) TABLET~\cite{slack2023tablet}, and (9) TabLLM~\cite{hegselmann2023tabllm}. In-context learning simply integrates a few training examples into the input prompt without model adaptation. TABLET improves on this by providing supplementary cues within the prompt, incorporating outputs from external classifiers in the form of rule sets and prototypes to better guide the inference process. TabLLM, in contrast, utilizes efficient parameter tuning methods such as IA3~\cite{liu2022few} to adapt the LLM using a limited set of labeled examples.

\vspace*{-0.12in}
\paragraph{Implementation details.}
\textsf{FeatLLM} employs GPT-3.5 as its underlying LLM, but its design remains flexible to accommodate alternative LLM architectures (see Appendix~\ref{sec:backbones} for performance using various backbones). For LLM inference, we configure the temperature parameter to 0.5 and set the top-p value to the API's default of 1. The ensemble size and the count of rules extracted are set at 20 and 10, respectively. Further analysis of the influence of hyper-parameters is shown in Figure~\ref{fig:hyperparameter2} and provided in Appendix~\ref{sec:hyper-parameter}. For training the linear model, we use the Adam optimizer with a learning rate of 0.01, and models are trained over 200 epochs. Epoch selection leverages $k$-fold cross-validation to achieve optimal model performance. For baseline comparisons, GPT-3.5 facilitates in-context learning, while T0~\cite{sanh2022multitask} serves as the backbone for TabLLM, in accordance with the settings in the original work. It is important to highlight that TabLLM is restricted to LLMs with open-source checkpoints, rendering GPT-3.5 incompatible for these experiments. When implementing conventional machine learning algorithms—including LogReg, XGBoost, and RandomForest—parameters are selected via Grid-search together with $k$-fold cross-validation. The value assigned to $k$ is either 2 or 4, ensuring that each training partition has representation from all classes. Additional technical and procedural details can be found in Appendix~\ref{sec:baselines}.

\vspace*{-0.12in}
\paragraph{Results.}
Tables~\ref{tab:main1} and \ref{tab:main2} display comparative results on various datasets. All experiments were conducted three times, each with different random splits of training and test data. The reported AUC metrics reflect the mean and standard deviation across runs. \textsf{FeatLLM} is either the highest or second highest performing model in most scenarios relative to the baselines considered. Remarkably, our framework matches the accuracy of standard tabular models using only 4 shots, whereas those baselines require 32 shots for similar performance (refer to Fig.~\ref{fig:compares}). While STUNT and TabPFN achieved noteworthy results on select datasets, these gains did not consistently translate into notable improvements over traditional machine learning approaches overall. Additionally, LLM-based methods such as in-context learning, TABLET, and TabLLM were unable to process tabular datasets with high dimensionality. In contrast, our approach—integrating both bagging and ensemble strategies—maintained solid performance even on feature-rich datasets. It should be highlighted that applying our bagging and ensemble methodology to other LLM-based methods is conceivable, but this would require doubling the number of inferences for each sample, dramatically increasing computational load to impractical levels.

\subsection{Analysis \& Discussion}
\label{sec:analysis}
\paragraph{Ablation study.} To assess the individual impact of core elements, we conducted an ablation analysis focused on: (1) \textsf{-Tuning}, where the linear model's weight tuning is excluded and logits are simply calculated by summing binary features; (2) \textsf{-Ensemble}, where the ensembling step is removed; (3) \textsf{-Description}, in which feature descriptions are left out; and (4) \textsf{-Reasoning}, where the first reasoning instruction step in rule generation is skipped.

Table~\ref{tab:ablation} provides the mean change in AUC resulting from each ablation, averaged across all datasets. Each ablated component, when altered, leads to reduced model performance. Notably, the contribution of weight tuning grows as the shot count rises. This underscores that a larger dataset enables more precise rule importance estimation, thus enhancing the effect of tuning. Conversely, the importance of both feature descriptions and reasoning instructions is more pronounced when only a few shots are available, underscoring the value of harnessing LLMs’ prior knowledge in low-data regimes. Finally, the ensemble strategy proposed here reliably enhances model outcomes under all tested conditions.

\vspace*{-0.12in}
\paragraph{Training \& inference time comparison.} We evaluate and contrast the durations required for training and inference across several methodologies. All tests utilize the Adult dataset with a training subset comprising 16 examples. For these experiments, a single A100 GPU serves as the default computational platform, with the exception of TabLLM, which is distributed across four A100 GPUs to enable model parallelism. The comparative results on runtime are presented in Table~\ref{tab:cost}. It is worth emphasizing that \textsf{FeatLLM} achieves an inference time that is relatively low and is on par with classic baseline methods. Conversely, LLM-driven methods such as In-context learning, TABLET, and TabLLM necessitate significantly more time for inference. With respect to training duration, \textsf{FeatLLM} performs similarly to other methods leveraging LLMs. Its requirement of just 30 API calls imposes minimal overhead and can be further optimized via parallelization.

\vspace*{-0.12in}
\paragraph{Quality of features generated by \textsf{FeatLLM}.}
To assess the practical utility of rule-based features produced by \textsf{FeatLLM}, we design a straightforward experiment. We compute the mean AUROC score that measures the association between each newly generated feature and the target labels, followed by calculating the proportion of features with an AUROC surpassing 0.5 for every dataset evaluated. 
A summary of these findings is presented in Table~\ref{tab:quality}. The results reveal that a substantial fraction of the generated rules are pertinent to the prediction target, thereby offering valuable information for solving the task (as observed in the “Before tuning” column). In addition, throughout linear model fine-tuning, any rules demonstrating minimal relevance—those with assigned weights below a chosen threshold—are automatically discarded (see the “After tuning” column). For these experiments, the threshold was specified as 0.1, determined by analyzing the overall weight distribution.

\vspace*{-0.12in}
\paragraph{Effect of adding spurious correlations.} To evaluate the capacity of different approaches to mitigate the impact of extraneous spurious correlations in a low-shot scenario, we performed a detailed investigation. Specifically, our experiment employed the Heart dataset while appending, at random, columns sourced from the Adult dataset that are unrelated to the Heart dataset's target variable. The experiment tracked model performance as the quantity of these unrelated columns increased incrementally. For consistency across methods, we fixed the number of labeled samples at 8 shots and did not include any unlabeled data, which meant excluding methods such as SCARF or STUNT from this analysis. 

Figure~\ref{fig:spurious} displays how model accuracy is affected as the presence of spurious correlations increases. In contrast to baseline methods, \textsf{FeatLLM} demonstrated superior robustness, exhibiting the smallest decline in performance. This result is presumably due to our framework’s explicit alignment of rule extraction with prior knowledge linking the features to the specific task, discouraging the creation of rules that depend on irrelevant columns and thereby improving reliability. In practice, rules involving noisy features constituted only about 1-2\% of all extracted rules. Nevertheless, we also notice that when columns from more semantically related sources are arbitrarily included, \textsf{FeatLLM} is susceptible to learning spurious relationships, potentially confusing causal structures between features and the target task (refer to Appendix~\ref{sec:spurious-appendix} for further discussion).

\vspace*{-0.12in}
\paragraph{Hyper-parameter analysis}
\textsf{FeatLLM} is governed primarily by two hyper-parameters: the ensemble size and the number of rules per class that are extracted with each LLM query. We systematically explored how variations in these parameters influence model efficacy. Our findings indicate that increasing the ensemble size generally enhances performance, but performance gains plateau past approximately 20 ensembles (refer to Fig.~\ref{fig:appendix-hyperparameter1} in the Appendix). Regarding the rule count, while differences did not reach statistical significance, we selected 10 rules as an optimal setting to balance prompt complexity and effectiveness (see Fig.~\ref{fig:hyperparameter2}). We attribute this to the trade-off where adding more rules increases input dimensionality—injecting more information—but simultaneously heightening the risk of overfitting, particularly in low-shot learning situations.

\section{Conclusion}
We introduce \textsf{FeatLLM}, which sets a new benchmark in few-shot tabular learning leveraging LLMs. In contrast to standard approaches where LLMs serve merely as sample-level predictors, \textsf{FeatLLM} positions them as feature engineers by generating classification criteria in the form of rules. Employing rule extraction and an ensemble strategy via bagging, \textsf{FeatLLM} remains lightweight in terms of inference resources, operates entirely via API calls with no need for further training, and sidesteps challenges posed by large feature sets. The approach delivers strong predictive accuracy, establishing \textsf{FeatLLM} as an efficient and practical solution for applying LLMs to tabular datasets with limited data availability.

Currently, \textsf{FeatLLM} is tailored for scenarios involving few available samples. Looking ahead, we aim to enhance the framework to facilitate automatic feature discovery for larger datasets, incorporate feature forms beyond simple rules, and improve interpretability, thus extending its versatility to a broader set of application domains.

\section{Broader Impact} The methodology presented in \textsf{FeatLLM} leverages the extensive prior knowledge embedded within LLMs to elevate tabular learning well past the boundaries of classical supervised methods. By focusing on maximizing learning efficiency, our work addresses key challenges related to overfitting when training data is scarce, drawing on supplemental prior knowledge for improved outcomes. This framework holds particular promise for practitioners contending with real-world datasets—such as those in finance and healthcare—where obtaining labeled examples is cost-prohibitive or otherwise difficult, yet relevant domain information is already incorporated within the pretraining of LLMs. An essential direction for future research concerns the examination of how inherent societal biases or inaccuracies present in LLMs’ knowledge bases might influence tabular model outputs, especially since these biases may have an outsized effect relative to available labeled data.

\end{document}